\documentclass[11pt]{article}
\usepackage[utf8]{inputenc}
\usepackage[spanish]{babel}
\usepackage[T1]{fontenc}
\usepackage[a4paper, margin=2.5cm]{geometry}
\usepackage{booktabs}
\usepackage{graphicx}
\usepackage{hyperref}
\usepackage{parskip}

\title{\textbf{UNIVERSIDAD POLITÉCNICA DE QUERÉTARO}\\[4pt]
\Large REPORTE ED.03.02: PRUEBA DE IMPLEMENTACIÓN DEL DASHBOARD DEL SIE DE BANXICO (SF43718)}
\author{Materia: Herramientas de Visualización --- IDIA223-4\\
Profesor: Jairo Pacindo Piña\\
Estudiante: Marco Alexander Banda Pérez --- Matrícula: 125054766\\
Fecha: 22 de noviembre de 2026}
\date{}

\begin{document}
\maketitle

\section{Ficha Resumida de la API}

\begin{table}[h]
\centering
\begin{tabular}{p{3.2cm}p{11.5cm}}
\toprule
\textbf{Campo} & \textbf{Detalle} \\
\midrule
API y documentación & SIE, Banco de México --- Sistema de Información Económica. \\
Endpoint & \texttt{https://www.banxico.org.mx/SieAPIRest/service/v1/series/\{claves\}/datos/\{inicio\}/\{fin\}} \\
Autenticación y límite & Token gratuito de 64 caracteres, pedido en \url{https://www.banxico.org.mx/SieAPIRest/service/v1/token}. Se envía en el encabezado \texttt{Bmx-Token}. El tablero valida que tenga al menos 32 caracteres y se detiene si falta. \\
Serie utilizada & \texttt{SF43718}, tipo de cambio FIX (pesos por dólar). Es la serie por defecto del tablero; el panel lateral permite pedir otras claves separadas por comas. \\
Rango & 1 de enero de 2021 a 31 de diciembre de 2025. \\
Variables utilizadas & \texttt{Fecha} (datetime), \texttt{Valor} (numérico, pesos por dólar), \texttt{Clave}, \texttt{Serie}, \texttt{Anio}, \texttt{Mes}. Derivadas: media móvil, variación diaria (\%), cambio mensual (\%), cambio anual (\%). \\
\bottomrule
\end{tabular}
\caption{Ficha de la API del SIE de Banxico, según el código del tablero Base.}
\end{table}

El tablero hace una sola petición con todas las claves, guarda el
resultado en caché por una hora y distingue los errores: 401 si el token
no fue aceptado, 429 si hay demasiadas peticiones, y error si el servidor
devuelve algo que no es JSON aunque el código HTTP sea 200.

\section{Documentación de Limpieza y Salida de Terminal}

El script valida el token antes de consultar, verifica los códigos HTTP
(401, 400, 429), confirma el encabezado \texttt{Content-Type:
application/json} y solo entonces procesa los datos. Al cargar los datos,
todos los valores llegan como texto, por lo que se convierten a formato
numérico; los días inhábiles llegan marcados como \texttt{N/E} y se
cuentan por serie en lugar de tratarse como ceros.

\begin{itemize}
\item \textbf{Tipos de dato originales:} \texttt{fecha}: texto,
  \texttt{dato}: texto (con comas como separador de miles y
  \texttt{N/E} en días inhábiles).
\item \textbf{Conversión de tipos:} se quitan las comas y
  \texttt{dato} se convierte a número; los \texttt{N/E} quedan como
  valores faltantes y se cuentan en \texttt{faltantes} por serie.
\item \textbf{Limpieza:} se eliminan las filas sin valor numérico
  antes de calcular indicadores y gráficas; la serie de referencia
  para los indicadores es la primera clave pedida.
\end{itemize}

Salida de terminal de la prueba de implementación realizada para esta
entrega, con datos reales del SIE (serie SF43718, 1 de enero de 2021 a
31 de diciembre de 2025):

\begin{verbatim}
$ python3 -m py_compile dashboard_banxico.py
SYNTAX_OK

registros brutos: 1258
N/E: 0
registros utilizables: 1258
rango: 2021-01-04 a 2025-12-31
primero: 19.8457 el 2021-01-04
ultimo: 18.0012 el 2025-12-31
cambio total periodo: -9.29 %
min: 16.3357 (2024-04-08)  max: 21.8185 (2021-11-26)
\end{verbatim}

La lógica de la gráfica nueva se verificó además con datos sintéticos
antes de la corrida real (cierres de año 20.00, 19.00, 17.00, 18.00,
16.50: cambios anuales -5.00\%, -10.53\%, +5.88\%, -8.33\%), para
confirmar que la fórmula estaba bien calculada; esos valores sintéticos
\textbf{no} son valores del tipo de cambio FIX y no se usan en las
lecturas de la sección 3.

\section{Análisis e Interpretación de las Cuatro Gráficas}

El tablero Base traía cuatro gráficas. En esta entrega se quitó el mapa
de calor ``Promedio por mes y año'' y se añadió ``Cambio anual (cierre
de año)'', por lo que el tablero final sigue teniendo cuatro gráficas.

\subsection{La serie y su tendencia}

\textbf{Pregunta:} ¿Cómo cambió el tipo de cambio FIX a lo largo del
tiempo y cuál fue su tendencia reciente?

\textbf{Lectura:} La línea gris es el valor diario y la línea azul es
la media móvil de 30 días, ajustable con el control lateral. La media
móvil quita el ruido diario y deja ver si el peso se estaba
debilitando o fortaleciendo en cada tramo. El FIX empezó en 19.8457
pesos por dólar el 4 de enero de 2021 y cerró en 18.0012 el 31 de
diciembre de 2025, una baja total de -9.29\% en el periodo. En el
camino tocó su máximo de 21.8185 el 26 de noviembre de 2021 y su
mínimo de 16.3357 el 8 de abril de 2024: el peso se fortaleció de
forma sostenida durante 2023, se debilitó con fuerza a mediados de
2024 hasta volver a cerca de 20.8 pesos, y volvió a fortalecerse
durante 2025.

\begin{figure}[h]
\centering
\includegraphics[width=0.95\textwidth]{figuras/figura_31_tendencia.png}
\caption{La serie y su tendencia (SF43718).}
\end{figure}

\subsection{Cambio de un mes a otro}

\textbf{Pregunta:} ¿En qué meses el tipo de cambio se movió más de un
mes a otro?

\textbf{Lectura:} Cada barra es el cambio porcentual del cierre de un
mes contra el cierre del mes anterior. Rojo significa que el valor
subió, es decir, que el peso se debilitó; azul significa que bajó, es
decir, que el peso se fortaleció. A diferencia de la tendencia, aquí
se ven los saltos mes a mes, no el nivel. El mes con mayor subida fue
junio de 2024, con +7.23\%, y el mes con mayor bajada fue diciembre
de 2021, con -4.56\%. Las subidas mensuales más fuertes se concentran
en 2024, lo que coincide con el salto que se ve en la tendencia de
ese año.

\begin{figure}[h]
\centering
\includegraphics[width=0.95\textwidth]{figuras/figura_32_cambio_mensual.png}
\caption{Cambio de un mes a otro (\%).}
\end{figure}

\subsection{Qué tan movido estuvo cada año}

\textbf{Pregunta:} ¿En qué años el tipo de cambio fue más volátil día
a día?

\textbf{Lectura:} Cada caja resume la distribución de la variación
diaria de ese año: la línea central es la mediana, la caja cubre el
rango típico y los puntos fuera son días atípicos. Un año con la caja
ancha y muchos puntos fuera fue un año movido, aunque su cambio anual
neto haya sido pequeño. Esta gráfica mide volatilidad, no dirección.
El año más volátil fue 2024, con la desviación estándar diaria más
alta (0.82\%) y el rango intercuartílico más ancho (0.95 puntos
porcentuales); también tuvo el día de mayor subida diaria del
periodo, +4.16\% el 7 de junio de 2024. El año menos volátil fue
2025 (desviación estándar diaria de 0.58\%). El día de mayor bajada
diaria fue el 3 de noviembre de 2023, con -2.89\%.

\begin{figure}[h]
\centering
\includegraphics[width=0.95\textwidth]{figuras/figura_33_boxplot_anual.png}
\caption{Boxplot de la variación diaria por año (\%).}
\end{figure}

\subsection{Cambio anual (cierre de año) --- gráfica añadida}

\textbf{Pregunta:} ¿En qué año se depreció o apreció más el peso en
el periodo?

\textbf{Lectura:} Cada barra es el cambio porcentual del cierre de
diciembre contra el cierre de diciembre del año anterior, con la
misma convención de color del tablero: rojo si el FIX subió en el
año (el peso se debilitó), azul si bajó (el peso se fortaleció).
Ninguna gráfica del Base respondía esto de forma directa: la
tendencia muestra el nivel pero no el porcentaje por año, el cambio
mensual tiene unas 60 barras y es ruidoso, y el boxplot muestra
volatilidad, no dirección ni magnitud anual. Por eso esta gráfica
sustituye al mapa de calor, que solo repetía el nivel en promedio
mensual, escondía la volatilidad dentro de cada mes y se leía mal
por color, sin valores exactos. Con los datos reales, 2024 fue el año
de mayor depreciación del peso, con +22.86\%, y 2025 el de mayor
apreciación, con -13.40\%, seguido de cerca por 2023 con -13.11\%.
Es decir, el peso se fortaleció en tres de los cuatro años con cambio
anual completo (2022, 2023 y 2025) y solo se debilitó de forma
fuerte en 2024, lo que explica que el balance total del periodo sea
una baja del FIX de -9.29\% a pesar de ese año.

\begin{table}[h]
\centering
\begin{tabular}{lrr}
\toprule
\textbf{Año} & \textbf{Cierre (MXN/USD)} & \textbf{Cambio anual (\%)} \\
\midrule
2021 & 20.4672 & --- \\
2022 & 19.4715 & -4.86 \\
2023 & 16.9190 & -13.11 \\
2024 & 20.7862 & +22.86 \\
2025 & 18.0012 & -13.40 \\
\bottomrule
\end{tabular}
\caption{Cambio anual del FIX, con datos del SIE (SF43718).}
\end{table}

\begin{figure}[h]
\centering
\includegraphics[width=0.95\textwidth]{figuras/figura_34_cambio_anual.png}
\caption{Cambio anual del FIX (\%). Gráfica añadida en ED.03.02.}
\end{figure}

\section{Reflexión sobre el Proceso de Desarrollo}

Lo más difícil de esta entrega fue no confundir lo que cada gráfica
mide. La tendencia y el mapa de calor parecían distintas, pero las dos
hablaban del nivel del tipo de cambio, así que el mapa de calor no
añadía información para decidir nada; por eso se quitó. En cambio,
el cambio anual sí respondía una pregunta abierta: cuánto se movió el
peso en cada año completo, con dirección y magnitud en un solo número
por año.

También fue necesario cuidar los detalles de implementación: el valor
del año se toma del cierre, no del promedio, igual que el cambio
mensual del tablero Base; y el cálculo anual se hizo agrupando por año
en lugar de usar un alias de frecuencia, para que funcione en
cualquier versión de pandas. El token nunca va en el código ni en el
repositorio: va en \texttt{.streamlit/secrets.toml} en local y en los
secretos de Streamlit Cloud al publicar.

\end{document}
